\section{Implications for Evaluation and Interpretation}
\label{sec:implications}

\subsection{Persistence must be measured alongside performance}
A regulatory framework should be evaluated across time, not only at a final task score. Useful outcomes include task quality, decision latency, resource expenditure, retention of prior competence, constraint violations, and recovery following a specified disturbance. Reporting these quantities separately prevents an aggregate score from hiding an operationally important failure.

For a known stationary model, distance to a stable equilibrium may help characterize recovery. In a dynamic task, that distance can be misleading: a useful action may move the system away from a frozen equilibrium, and a failing system may have a small net update because it is stuck. Equilibrium proximity is therefore a model-dependent diagnostic rather than a general intelligence metric.

\subsection{Accounting for the regulatory layer}
Computing information gain, estimating risk, monitoring reserves, and adapting weights all incur costs. These must be included in the same budget used to evaluate the agent. Otherwise an apparently resource-aware design can receive an unreported computational advantage. Model calibration and online operation should also be distinguished: expensive offline fitting does not have the same implications as an expensive real-time decision rule.

The five components should support diagnosis only where the mapping from observables to mechanisms is justified. A decline in performance cannot be attributed to the sustainability component merely because it occurs late in a deployment. Alternative explanations include distribution shift, memory corruption, and changing task difficulty. Controlled perturbations and held-out predictions are needed to separate them.

\subsection{Persistence is not a complete account of intelligence}
Stability and viability are relevant to sustained cognitive performance, but they do not explain all learning, reasoning, representation, or social understanding. Some intelligent actions deliberately leave a familiar operating regime. A system can also remain stable while pursuing harmful or irrelevant objectives. UVIF should therefore complement task- and context-specific accounts of cognition, rather than replace them with an unrestricted definition based on equilibrium.
